% =====================================================================================
% REVISIÓN 2  -  Propuesta de solución
% Escribe tu texto FUERA de los recuadros \begin{guia} ... \end{guia}.
% =====================================================================================

\section{Propuesta de solución} \label{sec:propuesta} % audit:req=propuesta_solucion

\begin{guia}[Guía: Propuesta de solución (Revisión 2)]
\textbf{¿Para qué sirve?} Explica \emph{qué vas a hacer exactamente y cómo resuelve el problema}, antes de entrar a los
detalles de construcción (Revisión 3).

\textbf{Debe responder:}
\begin{itemize}[nosep,leftmargin=*]
  \item ¿En qué consiste la solución, a grandes rasgos? Incluye un diagrama general (arquitectura, flujo del proceso o
        esquema del método).
  \item \emph{Proyecto de sistema:} ¿cuáles son los requisitos funcionales (qué hará) y no funcionales (rendimiento,
        seguridad, usabilidad)? ¿quiénes son los usuarios y qué roles tienen?
  \item \emph{Proyecto de investigación:} ¿cuál es la hipótesis o pregunta de investigación? ¿qué enfoque o algoritmo
        propones y en qué se diferencia de lo existente?
  \item ¿Qué tecnologías, herramientas y metodología de trabajo usarás y por qué (frente a las alternativas)?
  \item ¿Cómo atiende la propuesta cada objetivo específico? (una tabla ``objetivo $\rightarrow$ componente de la solución''
        es muy útil).
  \item ¿Cuál es el plan de trabajo y el cronograma del cuatrimestre?
  \item ¿Cómo sabrás que funcionó? (criterios de aceptación o métricas de éxito).
\end{itemize}

\textbf{Organización sugerida:} \verb|\subsection| para \emph{Descripción general}, \emph{Requisitos} (o \emph{Hipótesis y
enfoque}), \emph{Tecnologías y metodología} y \emph{Plan de trabajo}.

\textbf{Extensión mínima:} 150 palabras, además de diagramas y tablas.

\textbf{Errores frecuentes:} describir solo tecnologías sin explicar la solución; requisitos vagos (``el sistema será
rápido''); propuesta desconectada de los objetivos; cronograma sin fechas.

\textbf{Cómo se revisa:} coherencia con el problema y los objetivos, viabilidad en un cuatrimestre y claridad técnica.
\end{guia}

% >>> ESCRIBE AQUÍ la propuesta de solución (con una \subsection por apartado) <<<
